\section{Stream-Wrapper Evaluation Details}
\begin{table}[t]
\centering
\caption{Stream-wrapper authorization-order violations (\%), tenants T1, T8, T2, T5 and T4 (personalisation). Trials per schedule and AVD: pen 50, kill 50, fstop 50, reb 50, rev 50, net 50, infl 50.}
\label{tab:matrix}
\mtab
\begin{tabular}{@{}lrrrrrrr@{}}
\toprule
\mhead Variant & \textsf{pen} & \textsf{kill} & \textsf{fstop} & \textsf{reb} & \textsf{rev} & \textsf{net} & \textsf{infl} \\
\midrule
none & 100 & 100 & 80 & 100 & 100 & 80 & 44 \\
cancel & 20 & 20 & 20 & 20 & 20 & 80 & 1 \\
enqueue & 100 & 100 & 40 & 100 & 100 & 80 & 44 \\
job-start & 40 & 40 & 40 & 40 & 40 & 79 & 43 \\
conv & 20 & 20 & 20 & 20 & 20 & 80 & 2 \\
\midrule
egress (decl.) & 40 & 40 & 40 & 40 & 40 & 20 & \mgood{\textbf{0}} \\
\sys & \mgood{\textbf{0}} & \mgood{\textbf{0}} & \mgood{\textbf{0}} & \mgood{\textbf{0}} & \mgood{\textbf{0}} & \mgood{\textbf{0}} & \mgood{\textbf{0}} \\
\sys (AVD B) & \mgood{\textbf{0}} & \mgood{\textbf{0}} & \mgood{\textbf{0}} & \mgood{\textbf{0}} & \mgood{\textbf{0}} & \mgood{\textbf{0}} & \mgood{\textbf{0}} \\
\midrule
volatile labels & \mgood{\textbf{0}} & 100 & 80 & 100 & 100 & \mgood{\textbf{0}} & \mgood{\textbf{0}} \\
taint rule & \mgood{\textbf{0}} & \mgood{\textbf{0}} & \mgood{\textbf{0}} & \mgood{\textbf{0}} & \mgood{\textbf{0}} & \mgood{\textbf{0}} & \mgood{\textbf{0}} \\
\bottomrule
\end{tabular}
\end{table}


The stream wrapper measures \emph{authorization order} violations among delivered
items. Its event called \texttt{COMMIT} precedes the stream write. Accordingly, ``authorized''
in these stream-wrapper results does not establish a socket-handoff boundary. The channel tests
reported separately use accepted byte counts. We ask four questions.
\textbf{RQ1:} Which enforcement points authorize the upload of data that a withdrawal invalidated, under which lifecycle schedules?
\textbf{RQ2:} Do measured race windows match Theorem~\ref{thm:race}?
\textbf{RQ3:} Does \sys withhold permitted data, and what do its ablations lose?
\textbf{RQ4:} What does \sys cost, do its labels survive \texttt{SIGKILL}, and does it carry over to a real library?

\subsection{Authorization-order summary}
\label{app:stream}
The stream-wrapper \nTrialsValid{} trials compared nine variants over seven lifecycle schedules.
Under the authorization estimand, the conventional combination admits the mechanism-specific cases in
Table~\ref{tab:mechanisms}; the stream-wrapper \sys gate admits none on mediated paths. These observations motivate
lineage and late placement but cannot quantify physical handoffs in the log-to-write gap.
Permitted-flow controls distinguish suppression from non-delivery: stream-wrapper \sys logged
\rdFpBlocks{} policy blocks for the controls, while \deficitRdPct{}\% of permitted items
were undelivered on AVD~A (AVD~B: \deficitRdbPct{}\%). JobScheduler stopping uploads after reboot contributed
to that deficit. The synthetic contract-unaware policy lost \deficitTaintPct{}\%; this
is a policy ablation, not a comparison against a complete taint-tracking system.
\begin{table}[t]\centering\mtab
\caption{Conventional-check failures by mechanism (AVD~A): violating tenant-trials/total, stream-wrapper authorization estimand. The epoch row uses \vname{jobstart} outside \textsf{net}, because \vname{conv} purges the outbox of T5 on withdrawal and its T5 failures in \textsf{net} are retries. The selected fixture mix determines these rates.}\label{tab:mechanisms}
\begin{tabular}{@{}llr@{}}\toprule\mhead Mechanism & Variant, tenant, schedules & Trials\\\midrule
Lineage & \vname{conv}, T1, all & 252/350 \\
Epoch after re-grant & \vname{jobstart}, T5, all but net & 267/300 \\
Retry check & \vname{conv}, T8, net & 50/50 \\
\bottomrule\end{tabular}\end{table}

\paragraph{Supporting checks.} The stream-wrapper substrate tests include \subWorkManagerLogged{} WorkManager attempts (\subWorkManagerInvalid{} invalid) and \subAlarmManagerLogged{} AlarmManager attempts, with no prohibited authorizations for the fenced mediated paths. Across \crashKills{} process kills, atomic row labels leave \crashOrphA{} unlabelled rows, compared with \crashOrphB{} when data and labels are written separately. The authorization-window prediction agrees in \predAgree{} of \predN{} cases. Stream-wrapper per-submit medians of \eeNoneMs{} and \eeRdMs{}\,ms apply only to that wrapper; they are not the cost of the 4\,KiB channel transport measured below.


\subsection{RQ1: Which check points leak?}
\label{sec:eval-matrix}


\begin{figure}[t]
\centering
\begin{tikzpicture}
\begin{axis}[maturedbar,width=\linewidth,height=3.5cm,bar width=4.5pt,ymin=0,ylabel={items / tenant-trial},symbolic x coords={none,cancel,enq,start,conv,egress,RD,volat.,taint},xtick=data,x tick label style={font=\scriptsize},legend style={at={(0.5,1.0)},anchor=south,legend columns=2},enlarge x limits=0.07]
\addplot[fill=mBrick,draw=mInk!75] coordinates {(none,10.36) (cancel,0.20) (enq,8.88) (start,2.79) (conv,0.20) (egress,2.80) (RD,0.00) (volat.,8.21) (taint,0.00)};
\addlegendentry{lifecycle (pen, kill, fstop, reb, rev)}
\addplot[fill=mTeal,draw=mInk!75] coordinates {(none,6.42) (cancel,1.62) (enq,6.38) (start,6.12) (conv,1.63) (egress,1.28) (RD,0.00) (volat.,0.00) (taint,0.00)};
\addlegendentry{race (net, infl)}
\end{axis}
\end{tikzpicture}
\caption{Stream-wrapper prohibited authorizations per tenant-trial (T1, T8, T2, T5, T4), AVD~A and B pooled. \sys (RD) authorizes none.}
\Description{Grouped bars compare authorization-order failures of enforcement variants across lifecycle and retry schedules.}
\label{fig:bars}
\end{figure}


Table~\ref{tab:matrix} gives the violation rate per variant and schedule over the five tenants whose contracts cover pending
data, and Figure~\ref{fig:bars} gives the number of prohibited items. Across \nTrialsValid{} valid trials and
\nSubmissions{} logged submissions, four patterns stand out.

\emph{Enqueue-time checks are defeated by persistence.} \vname{none} and \vname{enqueue} authorized prohibited data in
\noneLifePct{}\% and \enqueueLifePct{}\% of tenant-trials under \textsf{pen}, \textsf{kill}, \textsf{reb} and \textsf{rev}.
The job had been queued before the withdrawal, and the system restored and ran it after process death or reboot.
Figure~\ref{fig:windows} shows how long the window stays open: a median of \winPendingMed{}\,s when pending and
\winRebootMed{}\,s across a reboot, because JobScheduler batching and boot time add to the job's own latency. Force-stop is
the exception. It removed the jobs, and the startup path that re-enqueued them re-ran the enqueue check, so \vname{enqueue}
leaked in only \enqueueFstopPct{}\% of \textsf{fstop} tenant-trials. The platform's lifecycle thus decides whether an
early check happens to be re-run.

\emph{Cancellation and job-start checks miss lineage and epochs.} \vname{conv} (cancel + enqueue + job start) leaked in
\convViolPct{}\% of tenant-trials overall. In the lifecycle schedules almost all of it came from T1, whose summary was
filed under \textit{diagnostics} (\convTOneLifePct{}\% of T1 tenant-trials, against \convOtherLifeK{} of \convOtherLifeN{}
for the other four tenants). Cancellation and checks keyed on the job's declared
scope never look at that job. \vname{job-start} and the serialized \vname{egress} check additionally leaked T5's
pre-withdrawal records after the re-grant (\vname{egress}: \egressTFivePct{}\% of T5 tenant-trials), because a boolean
flag cannot tell old data from new (Proposition~\ref{prop:nec}). \vname{egress} serializes its stream-wrapper gate but still admits data
whenever lineage or epochs matter.

\emph{Late authorizations defeat every check before the socket.} In the race schedules, the tenants other than T1
(T8, T2, T5, T4) still leaked under \vname{conv} in \convRacePct{}\% of tenant-trials. In \textsf{net}, the job had
passed its job-start check and prepared its batch before the withdrawal. It then retried that batch after connectivity
returned (\vname{job-start} \jobstartNetPct{}\%, \vname{conv} \convNetPct{}\%). Cancellation stopped the job only between
batches, so the batch being retried still went out (\vname{cancel} \cancelNetPct{}\%).

\emph{\sys authorized no prohibited item.} Over \rdTenantTrials{} tenant-trials on mediated paths, across both AVDs and all
seven schedules, \sys authorized \rdItems{} prohibited items. The second AVD reproduces this (row ``\sys (AVD B)''). The
\vname{volatile} ablation illustrates the consequence of losing labels and then failing open. It was clean while the process lived (\textsf{pen}:
\volatilePenPct{}\%), but after \texttt{kill}, reboot or force-stop its in-memory labels were gone and it failed open in
\volatileLifePct{}\% of tenant-trials.

\paragraph{Other deferred-work substrates.}
\label{sec:eval-substrate}
The main campaign uses JobScheduler. We reran \vname{none}, \vname{conv} and \sys with the same tenants on two other
substrates under \textsf{pen}, \textsf{kill}, \textsf{fstop} and \textsf{reb} (Table~\ref{tab:substrate}; at least
\subWorkManagerMinPerSched{} and \subAlarmManagerMinPerSched{} valid trials per schedule). WorkManager~2.9.1 keeps its own
database and reschedules work itself; for these builds we disabled the fixture's startup re-enqueue path, so every upload
after a lifecycle event comes from WorkManager. It behaved like JobScheduler, and its work also survived force-stop
(\vname{none} leaked in \subWorkManagerNoneForcestop{}\% of \textsf{fstop} tenant-trials). \vname{conv} again leaked only
T1's misfiled summary. With AlarmManager, the app arms one alarm per job and re-arms lost alarms from its outbox when its
process starts. Alarms do not survive force-stop or reboot. After force-stop the startup path re-armed some alarms but not
others. After a reboot no alarm app logged a process start within the observation window in any trial (in one manual
reboot outside the campaign its \texttt{BOOT\_COMPLETED} receiver did run, so this is timing we did not control), and nothing
re-armed the alarms. \vname{none} therefore leaked in \textsf{pen} and \textsf{kill}, in \subAlarmManagerNoneForcestop{}\%
of \textsf{fstop} and in no \textsf{reb} tenant-trials, and lost \subAlarmManagerNoneDefPct{}\% of the FP tenants' permitted
items to lost alarms, an availability loss rather than enforcement. \vname{conv} leaked T1's summary whenever its alarm fired. On both substrates \sys authorized no prohibited item
(\subWorkManagerRdItems{} over \subWorkManagerRdTT{} and \subAlarmManagerRdItems{} over \subAlarmManagerRdTT{} tenant-trials),
while the unmediated T9 leaked under \sys whenever its upload ran (\subWorkManagerRdTNinePct{}\% and
\subAlarmManagerRdTNinePct{}\% of its tenant-trials). The substrate decides whether deferred work survives a lifecycle
event, not which gate sees authorization.

\begin{table}[t]
\centering
\caption{Violation rate (\%) by deferred-work substrate (tenants T1, T8, T2, T5, T4). WorkManager and AlarmManager runs use their own rescheduling; the fixture's startup re-enqueue path is disabled for WorkManager.}
\label{tab:substrate}
\mtab
\begin{tabular}{@{}lrrrr@{}}
\toprule
\mhead Variant & \textsf{pen} & \textsf{kill} & \textsf{fstop} & \textsf{reb} \\
\midrule
\mgroup{5}{JobScheduler}
none & 100 & 100 & 80 & 100 \\
conv & 20 & 20 & 20 & 20 \\
\sys & \mgood{0} & \mgood{0} & \mgood{0} & \mgood{0} \\
\mgroup{5}{WorkManager}
none & 100 & 100 & 100 & 100 \\
conv & 20 & 20 & 20 & 20 \\
\sys & \mgood{0} & \mgood{0} & \mgood{0} & \mgood{0} \\
\mgroup{5}{AlarmManager}
none & 100 & 100 & 4 & 0 \\
conv & 20 & 20 & 5 & 0 \\
\sys & \mgood{0} & \mgood{0} & \mgood{0} & \mgood{0} \\
\bottomrule
\end{tabular}
\end{table}


\subsection{RQ2: Race windows follow Theorem~\ref{thm:race}}
\label{sec:eval-race}

\begin{figure}[t]
\centering
\begin{tikzpicture}
\begin{axis}[matured,width=\linewidth,height=3.4cm,xlabel={withdrawal offset after job start (s)},ylabel={P(violation)},xmin=0,xmax=1.6,ymin=-0.03,ymax=1.05,legend style={at={(0.5,1.03)},anchor=south,legend columns=2},transpose legend]
\addplot[mOchre,mark=triangle*,mark size=1.7pt] coordinates {(0.1,0.000) (0.3,0.042) (0.5,0.000) (0.7,0.000) (0.9,0.000) (1.1,0.000) (1.3,0.000) (1.5,0.000)};
\addlegendentry{cancel}
\addplot[mBrick,mark=square*,mark size=1.3pt] coordinates {(0.1,1.000) (0.3,1.000) (0.5,1.000) (0.7,0.059) (0.9,0.000) (1.1,0.000) (1.3,0.000) (1.5,0.000)};
\addlegendentry{job-start}
\addplot[mPlum,mark=diamond*,mark size=1.7pt] coordinates {(0.1,0.000) (0.3,0.000) (0.5,0.000) (0.7,0.059) (0.9,0.000) (1.1,0.000) (1.3,0.000) (1.5,0.000)};
\addlegendentry{conv}
\addplot[mTeal,mark=*,mark size=1.4pt,line width=1.2pt] coordinates {(0.1,0.000) (0.3,0.000) (0.5,0.000) (0.7,0.000) (0.9,0.000) (1.1,0.000) (1.3,0.000) (1.5,0.000)};
\addlegendentry{RD}
\addplot[mSlate,dashed,line width=0.9pt,no marks] coordinates {(0.00,1.000) (0.05,1.000) (0.10,1.000) (0.15,1.000) (0.20,1.000) (0.25,1.000) (0.30,1.000) (0.35,1.000) (0.40,1.000) (0.45,1.000) (0.50,1.000) (0.55,1.000) (0.60,1.000) (0.65,0.147) (0.70,0.000) (0.75,0.000) (0.80,0.000) (0.85,0.000) (0.90,0.000) (0.95,0.000) (1.00,0.000) (1.05,0.000) (1.10,0.000) (1.15,0.000) (1.20,0.000) (1.25,0.000) (1.30,0.000) (1.35,0.000) (1.40,0.000) (1.45,0.000) (1.50,0.000) (1.55,0.000) (1.60,0.000)};
\addlegendentry{Eq.~(2), measured $c_k$}
\end{axis}
\end{tikzpicture}
\caption{In-flight schedule: share of tenant-trials (T8, T2, T5) with a prohibited authorization, by withdrawal offset (200\,ms bins). The dashed line is Theorem~\ref{thm:race}'s prediction from the job-start variant's measured last-authorization offsets.}
\Description{Observed authorization-order outcomes are compared with predictions over withdrawal offsets.}
\label{fig:inflight}
\end{figure}

\begin{figure}[t]
\centering
\begin{tikzpicture}
\begin{axis}[matured,width=\linewidth,height=3.4cm,xlabel={enqueue-time check $\to$ first authorization (s)},ylabel={CDF},ymin=0,ymax=1.02,xmin=0,legend style={at={(0.98,0.03)},anchor=south east}]
\addplot[mInk,line width=1pt,no marks] coordinates {(17.5,0.007) (18.7,0.027) (20.1,0.047) (21.9,0.067) (22.5,0.087) (23.8,0.107) (30.5,0.127) (32.5,0.147) (32.8,0.167) (34.8,0.187) (35.5,0.207) (36.2,0.227) (39.1,0.247) (40.1,0.267) (40.7,0.287) (41.9,0.307) (42.1,0.327) (42.3,0.347) (42.4,0.367) (42.5,0.387) (42.6,0.407) (42.6,0.427) (42.6,0.447) (42.7,0.467) (42.7,0.487) (42.8,0.507) (42.8,0.527) (42.8,0.547) (42.9,0.567) (43.0,0.587) (43.0,0.607) (43.1,0.627) (43.1,0.647) (43.2,0.667) (43.3,0.687) (43.3,0.707) (43.3,0.727) (43.3,0.747) (43.4,0.767) (43.4,0.787) (43.6,0.807) (43.6,0.827) (43.6,0.847) (43.9,0.867) (43.9,0.887) (44.0,0.907) (54.9,0.927) (56.0,0.947) (61.5,0.967) (62.8,0.987) (63.4,1.000)};
\addlegendentry{\textsf{pen}}
\addplot[mBrick,dashed,line width=1pt,no marks] coordinates {(20.1,0.007) (20.6,0.027) (25.1,0.047) (25.4,0.067) (25.8,0.087) (27.9,0.107) (34.9,0.127) (39.3,0.147) (41.6,0.167) (42.0,0.187) (42.1,0.207) (42.1,0.227) (42.2,0.247) (42.2,0.267) (42.3,0.287) (42.4,0.307) (42.4,0.327) (42.5,0.347) (42.5,0.367) (42.6,0.387) (42.6,0.407) (42.7,0.427) (42.7,0.447) (42.8,0.467) (42.8,0.487) (42.8,0.507) (42.8,0.527) (42.9,0.547) (43.0,0.567) (43.0,0.587) (43.0,0.607) (43.0,0.627) (43.1,0.647) (43.1,0.667) (43.1,0.687) (43.2,0.707) (43.2,0.727) (43.3,0.747) (43.5,0.767) (43.6,0.787) (43.7,0.807) (43.8,0.827) (43.8,0.847) (44.3,0.867) (44.4,0.887) (44.6,0.907) (44.6,0.927) (44.7,0.947) (45.0,0.967) (49.1,0.987) (50.5,1.000)};
\addlegendentry{\textsf{kill}}
\addplot[mOchre,densely dotted,line width=1pt,no marks] coordinates {(16.5,0.007) (17.0,0.027) (17.1,0.047) (17.3,0.067) (17.4,0.087) (17.6,0.107) (17.7,0.127) (17.8,0.147) (17.9,0.167) (18.0,0.187) (18.2,0.207) (18.2,0.227) (18.3,0.247) (18.4,0.267) (18.4,0.287) (18.4,0.307) (18.5,0.327) (18.6,0.347) (18.6,0.367) (18.8,0.387) (18.8,0.407) (18.9,0.427) (18.9,0.447) (19.0,0.467) (19.0,0.487) (19.1,0.507) (19.2,0.527) (19.2,0.547) (19.3,0.567) (19.5,0.587) (19.5,0.607) (19.7,0.627) (19.8,0.647) (19.9,0.667) (20.0,0.687) (20.1,0.707) (20.5,0.727) (21.4,0.747) (38.5,0.767) (41.0,0.787) (42.6,0.807) (47.1,0.827) (47.8,0.847) (47.8,0.867) (48.0,0.887) (48.1,0.907) (48.2,0.927) (48.4,0.947) (48.6,0.967) (50.2,0.987) (51.2,1.000)};
\addlegendentry{\textsf{fstop}}
\addplot[mTeal,dashdotted,line width=1pt,no marks] coordinates {(27.9,0.007) (30.8,0.027) (31.2,0.047) (31.9,0.067) (32.8,0.087) (33.2,0.107) (34.1,0.127) (34.3,0.147) (34.8,0.167) (36.1,0.187) (37.1,0.207) (38.1,0.227) (38.5,0.247) (39.1,0.267) (39.3,0.287) (39.7,0.307) (40.4,0.327) (41.1,0.347) (41.5,0.367) (43.6,0.387) (45.7,0.407) (47.5,0.427) (48.7,0.447) (50.5,0.467) (50.9,0.487) (51.2,0.507) (52.7,0.527) (53.8,0.547) (56.4,0.567) (56.6,0.587) (57.2,0.607) (57.8,0.627) (58.2,0.647) (87.9,0.667) (88.3,0.687) (88.9,0.707) (89.0,0.727) (89.2,0.747) (89.5,0.767) (89.7,0.787) (89.9,0.807) (90.2,0.827) (90.5,0.847) (91.1,0.867) (91.6,0.887) (92.6,0.907) (94.5,0.927) (95.3,0.947) (96.0,0.967) (169.5,0.987) (184.0,1.000)};
\addlegendentry{\textsf{reb}}
\end{axis}
\end{tikzpicture}
\caption{Window an enqueue-time check leaves open: time from enqueue to the job's first authorization (\vname{none}, T2/T5/T6/T8 uploads), by schedule.}
\Description{Cumulative distributions show enqueue-to-authorization delays across lifecycle schedules.}
\label{fig:windows}
\end{figure}


In \textsf{infl}, the withdrawal lands at a uniform offset after the upload job starts. Theorem~\ref{thm:race} predicts that a
job-start check leaks exactly when the offset is below the job's last authorization offset $c_k$. We measured $c_k$ directly from the
\vname{job-start} variant's authorization log ($n{=}\nCmax{}$ jobs, median \cmaxMedianMs{}\,ms: three 4-record batches against a
300\,ms server). Per tenant-trial, the rule ``violation iff the offset is below that trial's own $c_k$'' predicted the observed outcome in
\predAgree{} of \predN{} cases, and the aggregate prediction $\Pr[c_k>u]$ matches the curve (Figure~\ref{fig:inflight}).
\vname{cancel} and \vname{conv} almost never leaked here (Table~\ref{tab:matrix}). Each batch is authorized at its start, and
cancellation stops the loop at the next batch boundary, so the only window left is the few milliseconds between the stop
check and the next authorization. Under \textsf{net} the same variants leaked in every tenant-trial, because the batch being
retried was authorized again \emph{after} the withdrawal. Where the window sits depends on the code path: cancellation
narrows one window and leaves another. \sys
stays at zero at every offset. We also count \emph{outstanding} submissions: authorized before the withdrawal, but with the HTTP exchange completing only
after the withdrawal was acknowledged. There were \outstandingRd{} such submissions for \sys and \outstandingSubs{} across all
variants in the race schedules. These stream-wrapper counts indicate authorization/response overlap, not confirmed outstanding bytes.
The sequence log and fixed response delay cannot establish receipt relative to acknowledgement.

To isolate serialization (A3), we ran a micro-benchmark. A submitter thread loops on single-item submissions while a
withdrawal lands at a random offset between 0 and 40\,ms (\S\ref{sec:eval-cost}, Table~\ref{tab:cost}). With the check and authorization in one
critical section, no round produced an authorization ordered after the withdrawal. With the check outside it (\vname{noser}),
\noserZeroPct{}\% of rounds did so with no added delay and \noserFivePct{}\% with a 5\,ms gap inserted between check and authorization.
In our wrapper the unserialized check runs before the connection is opened, so even without an added gap its window includes
connection setup to the host, a few milliseconds per submission. Submissions run back to back, so the withdrawal often lands
in such a window. Placing a check ``at egress'' is therefore not enough. It has to share the withdrawal's critical section.

\subsection{RQ3: Permitted flows and ablations}
\label{sec:eval-fp}
\label{sec:eval-boundary}

\begin{table}[t]
\centering
\caption{Permitted items \emph{not} delivered per tenant-trial (FP controls; T7 in \textsf{rev} only). Last column: share of all permitted items. \vname{none} enforces nothing, so its deficit is availability loss common to every variant.}
\label{tab:fp}
\mtab\setlength{\tabcolsep}{3pt}
\begin{tabular}{@{}lrrrrr@{}}
\toprule
\mhead Variant & T3 flush & T4 backup & T6 none & T7 OS & \% total \\
\midrule
none & 0.00 & 0.01 & 0.00 & 0.00 & 0.0 \\
cancel & 0.02 & 0.00 & 0.01 & 0.00 & 0.1 \\
enqueue & 0.00 & 0.00 & 0.00 & 0.00 & 0.0 \\
job-start & 0.02 & 0.08 & 0.00 & 0.00 & 0.3 \\
conv & 0.13 & 0.15 & 0.15 & 0.00 & 1.1 \\
egress (decl.) & 0.00 & 0.00 & 0.00 & 0.00 & 0.0 \\
\sys & 0.00 & 0.00 & 0.00 & 0.00 & 0.0 \\
\sys (AVD B) & 0.00 & 0.02 & 0.00 & 0.00 & 0.1 \\
volatile labels & 0.11 & 0.09 & 0.05 & 0.00 & 0.6 \\
taint rule & 9.54 & 10.61 & 0.22 & 12.00 & 55.6 \\
\bottomrule
\end{tabular}
\end{table}


Table~\ref{tab:fp} counts permitted items that were never delivered for the four FP-control tenants. These are the
flush-allowed opt-out (T3), the backup purpose of the restriction tenant (T4), the no-withdrawal control (T6), and
pre-revocation location history after \texttt{pm revoke} (T7). \vname{none} lost \deficitNonePct{}\% of these items and \sys
\deficitRdPct{}\% (AVD~B: \deficitRdbPct{}\%); \sys's fence logged \rdFpBlocks{} blocks for these tenants. The other non-\vname{taint} variants lost at most \deficitConvPct{}\%
(\vname{conv}). None of these losses is a blocked flow. All \defOtherTT{} affected tenant-trials were \textsf{reb} trials in
which JobScheduler stopped a running upload because connectivity dropped after boot (stop reason \emph{connectivity}). The fixture's jobs stop at
the next batch boundary without rescheduling, so the rest was still unsent when the window closed. They include T6, which has no
withdrawal and therefore cannot be blocked by any check. The contract-unaware \vname{taint} rule
blocked the flush, the backup and the OS-revocation history alike (\deficitTaintPct{}\%). This is the false-positive cost of
treating every withdrawal or revocation as ``block all uses of this source''. It is also why we refuse to count the
revocation residue in T7 as a violation.

\paragraph{Boundary.} T9 uploads through a raw socket that bypasses the wrapper. It leaked under every variant that lacks an
early check, \sys included (\rdUnmedPct{}\% of \sys tenant-trials). Variants with cancellation or job-start checks caught it
when the job had not yet started, so on an uncovered path an early check is better than nothing. This is exactly the
supported-path limit of Theorem~\ref{thm:sound}. T10 encrypts its body. Under \vname{none} and \vname{enqueue} the sink saw
\opaqueUnknownSubs{} post-withdrawal submissions it could not attribute to any record in \opaqueTrials{} tenant-trials. We
count these as \emph{unknown}, not clean. \sys authorized \rdOpaqueUnknown{} of them, because its decision is made before
encryption, on the device.

\subsection{RQ4: Cost, durability and a real library}
\label{sec:eval-cost}

\begin{table}[t]
\centering
\caption{Stream-wrapper cost, serialization and durability (on-device, AVD~B).}
\label{tab:cost}
\mtab
\begin{tabular}{@{}>{\raggedright\arraybackslash}p{.43\linewidth}>{\raggedleft\arraybackslash}p{.53\linewidth}@{}}
\toprule
\mgroup{2}{Overhead (median [bootstrap 95\% CI])}
Row insert, no label / label (\textmu{}s) & 641 [632, 651] / 652 [643, 661] \\
Critical section: log only (\textmu{}s) & 480 [476, 483] \\
\quad + declared check / + \sys rule (\textmu{}s) & 579 [573, 584] / 583 [576, 589] \\
Submission e2e: none / egress (ms) & 12.4 [11.1, 13.3] / 11.1 [10.3, 11.6] \\
Submission e2e: \sys (ms) & 13.1 [12.0, 14.0] \\
Job extras growth (bytes) & 52 \\
\mgroup{2}{Serialization: authorization after withdrawal}
\sys (serialized), gap 0\,ms & \mgood{0/300 (0\% [0, 1])} \\
\sys (serialized), gap 1\,ms & \mgood{0/300 (0\% [0, 1])} \\
\sys (serialized), gap 5\,ms & \mgood{0/300 (0\% [0, 1])} \\
\vname{noser}, gap 0\,ms & \mbad{117/300 (39\% [34, 45])} \\
\vname{noser}, gap 1\,ms & \mbad{98/300 (33\% [28, 38])} \\
\vname{noser}, gap 5\,ms & \mbad{139/300 (46\% [41, 52])} \\
\mgroup{2}{Durability under SIGKILL (100 kills)}
Label in same row: unlabelled rows & 0 \\
Label in separate store: unlabelled rows & 86 (86\% of kills) \\
\bottomrule
\end{tabular}
\end{table}


Table~\ref{tab:cost} reports on-device costs on the same emulator, measured with an otherwise idle host sink (medians with bootstrap intervals~\cite{efron1979bootstrap}). Persisting a
label in the same row adds \insDeltaUs{}\,\textmu{}s to a single-row insert (median, \insNoUs{} vs.\ \insLabUs{}\,\textmu{}s).
The fence's critical section (one grant lookup per scope, the authorization log write, no network) takes \fenceRdUs{}\,\textmu{}s
for a four-item batch, against \fenceNoneUs{}\,\textmu{}s for logging alone. End to end, a fenced submission takes
\eeRdMs{}\,ms against \eeNoneMs{}\,ms unfenced (medians, against a zero-delay local sink). The bootstrap intervals overlap
(Table~\ref{tab:cost}), so on this path the fence is within measurement noise of connection and HTTP costs. The label adds \extrasDelta{} bytes to a summary job's persisted
extras.

\paragraph{Durability.} Theorem~\ref{thm:sound} needs labels written atomically with data (A2). We killed a writer process
with \texttt{SIGKILL} \crashKills{} times at random moments. Writing the label in the same row left \crashOrphA{} unlabelled
rows. Writing the data row first and the label to \texttt{SharedPreferences} afterwards left \crashOrphB{} unlabelled rows,
i.e., at least one orphan in \crashOrphBPct{}\% of kills. Each orphan forces the fence to choose between failing open,
which is unsound, and failing closed, which blocks permitted data.

\subsection{Case study: ACRA}
\label{sec:eval-acra-historical}

In the preference-only ACRA~5.14.2 embedding~\cite{acra}, the startup processor
checked reporting preferences at enqueue, while an already scheduled sender service
sent approved reports without another preference check. We queued a crash offline, disabled
reporting, restored connectivity and observed for 45\,s. Stock ACRA delivered
\acraStockLeak{} reports (median \acraLatencyMed{}\,s after reconnection); the stream-wrapper fenced
sender delivered \acraFenceLeak{} and logged \acraFenceBlockLogged{} blocks. Allowed controls
delivered \acraStockCtl{} and \acraFenceCtl{} reports.

The scheduled job disappeared in all \acraJobGoneReb{} reboot trials. Reopening the app
at least 70\,s after the crash satisfied ACRA's report-age requirement and exercised its
startup preference check. Disabled stock/fenced reports delivered in \acraStockReb{} and
\acraFenceReb{} cases; enabled controls delivered in \acraStockRebCtl{} and
\acraFenceRebCtl{}.
